% 来源及取材范围见分题文件；此为整理者转录的正文，不是下载原件。
\begin{lemma}[10.119.9]
Let $R$ be a domain with fraction field $K$. Let $M$ be an $R$-submodule of $K^{\oplus r}$. Assume $R$ is local Noetherian of dimension $1$. For any nonzero $x\in R$ we have $\operatorname{length}_R(R/xR)<\infty$ and
\[\operatorname{length}_R(M/xM)\leq r\operatorname{length}_R(R/xR).\]
\end{lemma}
\begin{proof}
If $x$ is a unit then the result is true. Hence we may assume $x\in\mathfrak m$ the maximal ideal of $R$. Since $x$ is not zero and $R$ is a domain we have $\dim(R/xR)=0$, and hence $R/xR$ has finite length. Consider $M\subset K^{\oplus r}$ as in the lemma. We may assume that the elements of $M$ generate $K^{\oplus r}$ as a $K$-vector space after replacing $K^{\oplus r}$ by a smaller subspace if necessary.

Suppose first that $M$ is a finite $R$-module. In that case we can clear denominators and assume $M\subset R^{\oplus r}$. Since $M$ generates $K^{\oplus r}$ as a vectors space we see that $R^{\oplus r}/M$ has finite length. In particular there exists an integer $c\geq0$ such that $x^cR^{\oplus r}\subset M$. Note that $M\supset xM\supset x^2M\supset\ldots$ is a sequence of modules with successive quotients each isomorphic to $M/xM$. Hence we see that
\[n\operatorname{length}_R(M/xM)=\operatorname{length}_R(M/x^nM).\]
The same argument for $M=R^{\oplus r}$ shows that
\[n\operatorname{length}_R(R^{\oplus r}/xR^{\oplus r})=\operatorname{length}_R(R^{\oplus r}/x^nR^{\oplus r}).\]
By our choice of $c$ above we see that $x^nM$ is sandwiched between $x^nR^{\oplus r}$ and $x^{n+c}R^{\oplus r}$. This easily gives that
\[\begin{split}
r(n+c)\operatorname{length}_R(R/xR)&\geq n\operatorname{length}_R(M/xM)\\
&\geq r(n-c)\operatorname{length}_R(R/xR).
\end{split}\]
Hence in the finite case we actually get the result of the lemma with equality.

Suppose now that $M$ is not finite. Suppose that the length of $M/xM$ is $\geq k$ for some natural number $k$. Then we can find
\[0\subset N_0\subset N_1\subset N_2\subset\ldots N_k\subset M/xM\]
with $N_i\ne N_{i+1}$ for $i=0,\ldots,k-1$. Choose an element $m_i\in M$ whose congruence class mod $xM$ falls into $N_i$ but not into $N_{i-1}$ for $i=1,\ldots,k$. Consider the finite $R$-module $M'=Rm_1+\ldots+Rm_k\subset M$. Let $N'_i\subset M'/xM'$ be the inverse image of $N_i$. It is clear that $N'_i\ne N'_{i+1}$ by our choice of $m_i$. Hence we see that $\operatorname{length}_R(M'/xM')\geq k$. By the finite case we conclude $k\leq r\operatorname{length}_R(R/xR)$ as desired.
\end{proof}

\begin{lemma}[10.119.11]
Let $R$ be a domain with fraction field $K$. Let $M$ be an $R$-submodule of $K^{\oplus r}$. Assume $R$ is Noetherian of dimension $1$. For any nonzero $x\in R$ we have $\operatorname{length}_R(M/xM)<\infty$.
\end{lemma}
\begin{proof}
Since $R$ has dimension $1$ we see that $x$ is contained in finitely many primes $\mathfrak m_i$, $i=1,\ldots,n$, each maximal. Since $R$ is Noetherian we see that $R/xR$ is Artinian and $R/xR=\prod_{i=1,\ldots,n}(R/xR)_{\mathfrak m_i}$ by Proposition 10.60.7 and Lemma 10.53.6. Hence $M/xM$ similarly decomposes as the product $M/xM=\prod(M/xM)_{\mathfrak m_i}$ of its localizations at the $\mathfrak m_i$.

By Lemma 10.119.9 applied to $M_{\mathfrak m_i}$ over $R_{\mathfrak m_i}$ we see each $M_{\mathfrak m_i}/xM_{\mathfrak m_i}=(M/xM)_{\mathfrak m_i}$ has finite length over $R_{\mathfrak m_i}$. Thus $M/xM$ has finite length over $R$ as the above implies $M/xM$ has a finite filtration by $R$-submodules whose successive quotients are isomorphic to the residue fields $\kappa(\mathfrak m_i)$.
\end{proof}

\begin{lemma}[10.119.12, Krull--Akizuki]
Let $R$ be a domain with fraction field $K$. Let $L/K$ be a finite extension of fields. Assume $R$ is Noetherian and $\dim(R)=1$. In this case any ring $A$ with $R\subset A\subset L$ is Noetherian.
\end{lemma}
\begin{proof}
Let $I\subset A$ be a nonzero ideal. By Lemma 10.30.8 we can find a nonzero element $x\in I\cap R$. Then we get $I/xA\subset A/xA$. By Lemma 10.119.11 the $R$-module $A/xA$ has finite length as an $R$-module. Hence $I/xA$ has finite length as an $R$-module. Hence $I$ is finitely generated as an ideal in $A$.
\end{proof}
